\section{Open questions}

Five questions that this replication is genuinely unable to close, each with a concrete
next-step probe. These are gaps that remain after the finite-size qualitative replication,
not rhetorical framing.

\begin{enumerate}

\item \textbf{Robustness of the transition to correlated / non-Markovian noise.}\\
\emph{Basis.} The paper works in the depolarizing / global-white-noise limit, and our
replication inherits that. Real Sycamore-class hardware has substantial coherent, correlated,
and non-Markovian noise components that could smear the transition or shift the analytic
$(\varepsilon N)_c = \ln(5/2)$ location.\\
\emph{Next steps.} Repeat the 1D brickwork sweep with (a) two-qubit correlated depolarizing
noise, (b) coherent over-rotation noise, and (c) 1/$f$-like slow drift. Locate the onset of
$\chi/F$ blowup in each case at $N=10,12$ and compare onset location and effective transition
width against the pure-depolarizing baseline.

\item \textbf{Interaction with modern classical spoofing algorithms.}\\
\emph{Basis.} The paper's transition marks the boundary between (i) the regime where $\chi$
is a fidelity proxy (hard to spoof classically without also achieving $F$) and (ii) the regime
where $\chi$ decouples from $F$ (potentially spoofable at low $F$). Whether known spoofers
(Pan--Chen--Zhang tensor-network, Gao--Kalai stat-mech) actually operate on the low-$F$ side
of the transition is an operational question the paper does not close.\\
\emph{Next steps.} At fixed circuit family, evaluate (i) the noisy-quantum $\chi$ from our
density-matrix pipeline and (ii) $\chi$ from a shallow tensor-network spoofer at matched
compute budget, both as functions of $\varepsilon N$. Test whether the spoofer's $\chi$ tracks
the noisy-quantum $\chi$ above the transition even when its $F$ is essentially zero.

\item \textbf{Sycamore-scale application.}\\
\emph{Basis.} Our replication caps at $N=10$ exact statevector; the paper's headline
$(\varepsilon N)_c = \ln(5/2)$ and $-0.92$/layer saturation rely on $N=40$. Whether real
Sycamore-scale demonstrations sit safely below or dangerously near the transition threshold
is the practical question motivating the paper.\\
\emph{Next steps.} Use MPS or gauge-Pauli sampling to reach $N=40$--$60$; compute $\chi$ and
$F$ for reported Sycamore-70 circuit families at their published per-cycle error rates; place
each demonstration on the $(\varepsilon N,\chi/F)$ diagram and report the distance in
$\varepsilon N$ from the theoretical transition.

\item \textbf{Finite-size scaling of the transition width.}\\
\emph{Basis.} The paper claims \emph{sharpness} (transition width shrinking as $N\to\infty$),
but our $N=4,6,8,10$ data looks like a smooth crossover. Without a finite-size-scaling analysis
we cannot distinguish a genuine phase transition from a crossover that only sharpens beyond
CPU reach.\\
\emph{Next steps.} Fit $\chi/F$ vs $(\varepsilon N - (\varepsilon N)_c)\,N^{\alpha}$ to a common
scaling function at $N=6,8,10,12$ (adding $N=12$ via MPS if statevector budget is too tight);
estimate the width exponent $\alpha$; check whether $\alpha$ is stable as $N$ grows or drifts
(indicating we are still in a pre-asymptotic regime).

\item \textbf{Connection to average-case classical sampling complexity.}\\
\emph{Basis.} The paper's result is a statement about the XEB \emph{estimator}, not directly
about sampling complexity. If $\chi$ stops tracking $F$ above the transition, the operational
meaning of an XEB-based quantum-advantage claim in that regime is unclear, but the underlying
sampling problem may or may not become classically easy at the same threshold.\\
\emph{Next steps.} Combine our $\chi/F$ sweep with runtime measurements of a state-of-the-art
tensor-network sampler at the same $(N,d,\varepsilon)$ points; test whether classical sampling
cost drops sharply exactly at $(\varepsilon N)_c$, or whether the two thresholds
(XEB--fidelity decoupling vs classical-easiness) are distinct. Report both thresholds in the
same $(\varepsilon N, N)$ plane.

\end{enumerate}
